\documentclass[10pt,a4paper]{amsart}
\usepackage[margin=19mm]{geometry}
\usepackage{amsmath,amssymb,mathtools}
\usepackage[scr=rsfs,cal=euler]{mathalfa}
\usepackage{xfrac}
\usepackage[unicode,hidelinks,bookmarks=false]{hyperref}
\newcommand{\ie}{i.e.\ }
\newcommand{\cf}{cf.\ }
\newcommand{\resp}{respectively\ }
\newcommand{\Subsection}{\subsection}
\providecommand{\nobreakdash}{\nobreak\hbox{-}}
\setlength{\emergencystretch}{2em}
\multlinegap0pt
\usepackage{amssymb}
\newtheorem{proposition}{Proposition}
 \renewcommand{\Im} {\mathrm{Im\,}}
\newcommand{\ep}{\varepsilon}
\title{$L^2$ restriction bounds for analytic continuations of quantum ergodic Laplace eigenfunctions}
\author{John A. Toth, Xiao Xiao}
\date{Source: arXiv:2510.05570v2 (CC BY 4.0)}
\begin{document}
\maketitle

\noindent\emph{Source and licence.} $L^2$ restriction bounds for analytic continuations of quantum ergodic Laplace eigenfunctions. Version arXiv:2510.05570v2 (CC BY 4.0), distributed by arXiv under the Creative Commons Attribution 4.0 International licence, http://creativecommons.org/licenses/by/4.0/, as recorded in the arXiv licence metadata for this exact version and shown on the abstract page https://arxiv.org/abs/2510.05570v2. Full licence notice: \texttt{licenses/arxiv-2510.05570-CC-BY-4.0.txt}.\par\smallskip

\noindent\textit{Editorial scope.} Counted unit: the article's proposition multiplier (source0.tex lines 1044--1048). The article's complete original proof of it is delivered with it (source0.tex lines 1050--1116, one \texttt{proof} environment of 67 lines). No mathematics is written, repaired or re-derived: the statement and the proof are the article's own lines, in the article's own notation, and no retained line is substituted or re-flowed.  Retained uncounted as the source-local context the proof needs: lines 739--747; lines 841--845; lines 870--872; lines 890--896; lines 933--936.  Nothing was omitted from any retained span, so the package carries no omission record.  No retained line contains a citation, so the wrapper carries no bibliography and no bibliography entry.  No retained line contains a figure environment, an image-inclusion command or drawing code, so no image is delivered and the package declares no asset dependency.  Editorial formatting only, in addition: an AMS \texttt{amsart} preamble that restates the article's own declarations and macros used by the retained lines, namely the environments \texttt{proposition} on separate counters, together with the standard packages those lines need.  The retained environments print in this document as Proposition 1 in order of appearance, whereas the article prints the same items with the numbering of its own full document.  The complete native source of the article as distributed in the arXiv e-print for this exact version is bundled unchanged as \texttt{sources/186541/source0.tex}.
\begin{proposition}\label{WF1}
Let $(M,g)$ be a compact $C^{\omega}$ Riemannian manifold and $\{ u_h \}$ be {\em any} sequence of $L^2$-normalized Laplace eigenfunctions on $M.$ Then, the semiclassical wavefront set of $Tu_h$ satisfies
    \begin{equation}\label{wf1}
        \begin{split}
            WF_h(Tu_h)\subset\big\{(x,\xi :&x^*,\xi^*)\in T^*(B_{\tau}^*M): |\xi|_g=1, \, \xi^* = 0, \, x^* = \xi \big\}.
            \end{split}
    \end{equation}
 In particular, $WF_h(Tu_h)$ is a  compact subset of $T^* (B_{\tau}^*M).$
\end{proposition}

\begin{align} \label{wf4}
h^{2n} \chi_{\epsilon}^+( hD_{\beta_{\xi}} ) Tu(\alpha) \hspace{4in} \nonumber \\
= \int \cdots \int e^{i  \big( \langle \alpha - \beta, \beta^* \rangle  + \Phi(\beta,\tilde{\beta},x)  \big)/h } \, \chi_{\epsilon}^+ (\beta_{\xi}^*)  \chi_{\ep/2}^+ (x-\beta_x ) c Tu_h(\tilde{\beta})d\beta d\tilde{\beta} dx d\beta^* \nonumber \\
+ \int \cdots \int e^{i  \big( \langle \alpha - \beta, \beta^* \rangle  + \Phi(\beta,\tilde{\beta},x)  \big)/h } \, \chi_{\epsilon}^+ (\beta_{\xi}^*)  \chi_{\ep/2}(x-\beta_x ) cTu_h(\tilde{\beta})d\beta d\tilde{\beta} dx d\beta^*
\end{align}

\begin{equation} \label{W}
{\mathcal W}:= \big\{(x,\xi :x^*,\xi^*)\in T^*(B_{\tau}^*M): |\xi|_g=1, \, \xi^* = 0, \, x^* = \xi \big\}.
\end{equation}

\begin{proposition} \label{WF2}
Let $(M,g)$ be a compact $C^{\omega}$ Riemannian manifold, $\{ u_h \}$ be {\em any} sequence of $L^2$-normalized Laplace eigenfunctions on $M$ and $\Sigma=\{ F = 0 \}$ with $dF |_{\Sigma} \neq 0,$ be a compact hypersurface in the Grauert tube $B_{\tau}^*M.$ Then,    there exists a constant $T_F>0$ such that  \begin{equation}\label{wf1'}
        WF_h(T_{\Sigma}u_h)\subset \bigcup_{|\tau| \leq T_F} \exp \tau X_{F} \, \pi_{\Sigma}({\mathcal W}),
        \end{equation}
    where ${\mathcal W}$ is defined in (\ref{W}).
In particular, $WF_h(T_{\Sigma}u_h)$ is a  compact subset of $T^*\Sigma$.
\end{proposition}

\begin{eqnarray} \label{crit}
 - \beta_{x'}^* + \beta_{\xi'}  = O(|\beta_{x} -x|),  \,\,\,\, -\beta_{x_n}^* + G(\beta_x,\beta_{\xi'}) = O(|\beta_x -x|), \,\,\,\,
  \beta_{\xi'}^* = O(|\beta_x -x|). 
\end{eqnarray}

\begin{proposition} \label{multiplier}
Let $\Sigma \subset B_{\tau}^*M$ be a real, oriented,  compact hypersurface. Then, given $Q(h) \in \Psi_h^{0}(\Sigma),$ there exists a {\em function} $q_{\Sigma} \in C^{\infty}(\Sigma)$ such that
$$ Q(h) T_{\Sigma} u_h = q_{\Sigma} \,T_{\Sigma} u_h + o(1) \| T_{\Sigma} u_h \|_{L^2},$$
with $q_{\Sigma} = \sigma(Q) |_{{\mathcal W}_{\Sigma}'}.$ 
\end{proposition}

\begin{proof}
Following the argument in Proposition \ref{WF2}, we
assume first that $\partial_{\xi_n} F (x,\xi) \neq 0$ for all $(x,\xi) \in \Sigma \cap U$ near $q_0.$ Then, by the implicit function theorem, $ \Sigma \cap U = \{ (x, \xi'); \xi_n = G(x,\xi') \}$ where $G \in C^{\infty}_{loc}.$ Consequently, one can use $\alpha_{\Sigma} := (\alpha_x, \alpha_{\xi'})$ as local coordinates on $\Sigma \cap U$ and we denote canonical dual coordinates by $\alpha_{\Sigma}^*:= (\alpha_x^*,\alpha_{\xi'}^*).$ It follows that



\begin{align} \label{mult1}
Q(h) T_{\Sigma} u_h (\alpha_{\Sigma})  \hspace{4in} \nonumber \\
= \int e^{i \Psi_{\Sigma}(\alpha_{\Sigma}, \beta_{\Sigma}, \beta_{\Sigma}^*, \tilde{\beta},x)/h} \, q(\alpha_{\Sigma}, \beta_{\Sigma}^*) a (\beta_{\Sigma}, \tilde{\beta},x) Tu_h(\tilde{\beta}) d\tilde{\beta}  dx d\beta_{\Sigma} d \beta_{\Sigma}^*  + O(e^{-C/h}),
\end{align}
where the total phase
\begin{equation} \label{mult2}
\Psi_{\Sigma}(\alpha_{\Sigma}, \beta_{\Sigma},\beta_{\Sigma}^*, \tilde{\beta},x) = \langle \alpha_{\Sigma} - \beta_{\Sigma}, \beta_{\Sigma}^* \rangle + \Phi( \beta_\Sigma, G(\beta_{\Sigma}); \tilde{\beta},x).
\end{equation}

In view of the critical point equations in (\ref{crit}), one Taylor expands the symbol $q(\alpha_{\Sigma}, \beta_{\Sigma}^*)$ around $\beta_{x'}^* = \beta_{\xi'}, \, \beta_{x_n}^* = G(\beta_x, \beta_{\xi'})$ and $\beta_{\xi'}^* = 0.$ Setting $\beta_{\Sigma}^{0} = (  \beta_{\xi'},  G(\beta_x, \beta_{\xi'}), 0)$ and writing

$$ 
q(\alpha_{\Sigma}, \beta_{\Sigma}^*) = q(\alpha_{\Sigma}, \beta_{\Sigma}^{0}) + A ( \beta^{*}_{\Sigma} - \beta_{\Sigma}^{0}  ), \quad A = (A_1,A_2,A_3) \text{ with} \, A_j \in S^{0}(\Sigma); j=1,2,3,
$$

one writes the integral on the RHS of (\ref{mult1}) in the form


\begin{align} \label{mult3}
 \int e^{i \Psi_{\Sigma}(\alpha_{\Sigma}, \beta_{\Sigma}, \beta_{\Sigma}^*, \tilde{\beta},x)/h} \, q(\alpha_{\Sigma}, \beta_{\Sigma}^{0}) \, a (\beta_{\Sigma}, \tilde{\beta},x) Tu_h(\tilde{\beta}) d\tilde{\beta}  dx d\beta_{\Sigma} d \beta_{\Sigma}^* \nonumber \\
 + \int e^{i \Psi_{\Sigma}(\alpha_{\Sigma}, \beta_{\Sigma}, \beta_{\Sigma}^*, \tilde{\beta},x)/h} \, A \cdot ( \beta^{*}_{\Sigma} - \beta_{\Sigma}^{0}  )  \, a (\beta_{\Sigma}, \tilde{\beta},x) Tu_h(\tilde{\beta}) d\tilde{\beta}  dx d\beta_{\Sigma} d \beta_{\Sigma}^*=: I_1 + I_2
\end{align}


For the first term $I_1$ we make a standard Taylor expansion of the amplitude around $\alpha_{\Sigma} = \beta_{\Sigma}$ and integrate by parts with respect to $hD_{\beta_{\Sigma}^*}$ to get that

\begin{eqnarray} \label{I1}
I_1 =   \int e^{i \Psi_{\Sigma}(\alpha_{\Sigma}, \beta_{\Sigma}, \beta_{\Sigma}^*, \tilde{\beta},x)/h} \, q(\alpha_{\Sigma}, \alpha_{\Sigma}^{0}) \, a (\beta_{\Sigma}, \tilde{\beta},x) Tu_h(\tilde{\beta}) d\tilde{\beta}  dx d\beta_{\Sigma} d \beta_{\Sigma}^* + O(h) \| T_{\Sigma} u \|_{L^2} \nonumber \\
=  q(\alpha_{\Sigma}, \alpha_{\Sigma}^{0}) T_{\Sigma}u +  O(h) \| T_{\Sigma} u \|_{L^2}. \hspace{2in}
 \end{eqnarray}

 
 To estimate $I_2$ we make a further decomposition and write
\begin{eqnarray} \label{I2}
I_2 = \int e^{i \Psi_{\Sigma}(\alpha_{\Sigma}, \beta_{\Sigma}, \beta_{\Sigma}^*, \tilde{\beta},x)/h} \, \langle A, \beta^{*}_{\Sigma} - \beta_{\Sigma}^{0}  \rangle  \, \chi_{\epsilon} ( \beta^{*}_{\Sigma} - \beta_{\Sigma}^{0}  ) \, a (\beta_{\Sigma}, \tilde{\beta},x) Tu_h(\tilde{\beta}) d\tilde{\beta}  dx d\beta_{\Sigma} d \beta_{\Sigma}^* \nonumber \\
+ \int e^{i \Psi_{\Sigma}(\alpha_{\Sigma}, \beta_{\Sigma}, \beta_{\Sigma}^*, \tilde{\beta},x)/h} \, \langle A,   \,  \beta^{*}_{\Sigma} - \beta_{\Sigma}^{0}  \rangle  \, \chi_{\epsilon}^+ ( \beta^{*}_{\Sigma} - \beta_{\Sigma}^{0}  ) \, a (\beta_{\Sigma}, \tilde{\beta},x) Tu_h(\tilde{\beta}) d\tilde{\beta}  dx d\beta_{\Sigma} d \beta_{\Sigma}^* .
\end{eqnarray}

Since 
 $$\langle A, \beta^{*}_{\Sigma} - \beta_{\Sigma}^{0}  \rangle  \, \chi_{\epsilon} ( \beta^{*}_{\Sigma} - \beta_{\Sigma}^{0}  ) = O(\epsilon),$$
 it follows by $L^2$-boundedness that the first term
 \begin{eqnarray} \label{I2.1}
 \Big\|  \int e^{i \Psi_{\Sigma}(\cdot, \beta_{\Sigma}, \beta_{\Sigma}^*, \tilde{\beta},x)/h} \, \langle A, \beta^{*}_{\Sigma} - \beta_{\Sigma}^{0}  \rangle  \, \chi_{\epsilon} ( \beta^{*}_{\Sigma} - \beta_{\Sigma}^{0}  ) \, a (\beta_{\Sigma}, \tilde{\beta},x) Tu_h(\tilde{\beta}) d\tilde{\beta}  dx d\beta_{\Sigma} d \beta_{\Sigma}^* \Big\|_{L^2(U)} \nonumber \\
 = O(\epsilon) \|T_{\Sigma} u \|_{L^2}.\end{eqnarray}
 

 As for the second term in (\ref{I2}), noting that $\Im \phi(\beta,x) \geq \frac{1}{C} |\beta_x - x|^2,$ it follows by a similar integration by parts argument as in the proof of Proposition \ref{WF1} (see (\ref{wf4}) ) that 
 \begin{eqnarray} \label{I2.2}
\Big\| \int e^{i \Psi_{\Sigma}(\cdot, \beta_{\Sigma}, \beta_{\Sigma}^*, \tilde{\beta},x)/h} \, \langle A,   \,  \beta^{*}_{\Sigma} - \beta_{\Sigma}^{0}  \rangle  \, \chi_{\epsilon}^+ ( \beta^{*}_{\Sigma} - \beta_{\Sigma}^{0}  ) \, a (\beta_{\Sigma}, \tilde{\beta},x) Tu_h(\tilde{\beta}) d\tilde{\beta}  dx d\beta_{\Sigma} d \beta_{\Sigma}^* \Big\|_{L^2(U)}  \nonumber \\
 = O_{\epsilon}(h^{\infty}). \hspace{1in}
 \end{eqnarray}
 Consequently, it follows from (\ref{I2.1}), (\ref{I2.2}) and (\ref{I1}) that
 \begin{equation} \label{upshot}
\|  Q(h) T_{\Sigma} u_h  - q(\alpha_{\Sigma}, \alpha_{\Sigma}^0) T_{\Sigma} u_h \|_{L^2(U)} = ( O(\epsilon) + O_{\epsilon}(h^{\infty}) ) \| T_{\Sigma} u \|_{L^2}.
\end{equation}
Since $\epsilon >0$ is arbitrary, this proves the Proposition in the first case where locally $ \Sigma \cap U = \{ (x, \xi'); \xi_n = G(x,\xi') \}.$

In the second case where $ \Sigma \cap V = \{ (x',\xi); x_n = H(x',\xi)  \},$ one argues similarly to the above but with $\beta_{\Sigma}^0 = ( \beta_{\xi'} + \beta_{\xi_n} \partial_{\beta_x'}H, \beta_{\xi_n} \partial_{\beta_{\xi}}H ).$
By constructing  a partition of unity, one can cover $\Sigma$ by finitely many sets of the form $U$ and $V$. This completes the proof.

\end{proof}

\end{document}
